\documentclass[xcolor=svgnames]{beamer}
\usepackage{latexsym,amssymb,amsmath,stmaryrd}
\usepackage[utf8]{inputenc}
\usepackage[T1]{fontenc}
\usepackage[francais]{babel}
\usepackage{pgf}
\usepackage{tikz}
\usetikzlibrary{arrows,patterns,plotmarks,shapes,snakes,er,3d,automata,%
backgrounds,topaths,trees,petri,mindmap}
\usepackage{pgfbaseimage}
\usepackage{ifpdf}
\ifpdf
\usepackage{graphicx}
\else
%\usepackage[dvips]{graphicx}
\usepackage{pstricks,pst-tree,pst-node}
\newpsobject{showgrid}{psgrid}{subgriddiv=1,griddots=10,gridlabels=6pt}
\fi
\usepackage{moreverb}
\usepackage[normalem]{ulem}
\usepackage{version}
\usepackage{url}
\usepackage{multicol}
\usepackage{listings}
\lstset{
  language=[ANSI]C, 
%  gobble=2, 
%  escapeinside="", 
  basicstyle=\ttfamily, 
%  directivestyle=\color{Fuchsia},
  identifierstyle = \color{DarkOrange},
  keywordstyle=\color{DarkGreen}, 
  commentstyle=\color{red}, 
%  numbers=left, 
%  numbersep=5pt, 
%  numberstyle=\scriptsize, 
%  moredelim=[is][\color{gray}\itshape]{/*}{*/}, 
%  moredelim=[is][\alert]{/+}{+/}, 
%  morecomment=[is]{/=}{=/}, 
%  morecomment=[in]{comment=}{,} 
%  fancyvrb=true 
} 

\newcommand{\binaire}[1]{\ensuremath{\underline{#1}}}
\newcommand{\C}[1]{\texttt{#1}}
\newcommand{\bbbn}{\ensuremath{\mathbb{N}}}
\input{communs/cc_beamer}
\usetheme{classic}
\newcommand{\nowrite}{\put(10,-4){\includegraphics[scale=.05]{creative_commons/nopencil}}}
\newcommand{\youwrite}{\put(10,-4){\includegraphics[scale=.05]{creative_commons/pencil}}}
\newcommand{\aemporter}{\put(10,-6){\includegraphics[scale=.05]{creative_commons/szymonraj_Shopping_bag}}}


%%% Titre 
\title{Éléments d'informatique -- Cours 8. Types de données, représentations et entrées/sorties.}
\author{Pierre Boudes}
\date{\today}

\begin{document}

%% Page de titre et licence CC.
\begin{frame}
	\titlepage
	\vfill
	\begin{center}
		\CcGroupByNcSa{0.83}{0.95ex}\\[2.5ex]
		{\tiny\CcNote{\CcLongnameByNcSa}}
		\vspace*{-2.5ex}
	\end{center}
\end{frame}

\section[Plan]{}
\frame[label=plan]{\tableofcontents}

%%%%%%%%%%%%%%%%%%%%
 \section*{Contenu du cours}
\begin{frame}
%  \frametitle{Contenu du cours (rappel)}
 \begin{itemize}
  \item {Éléments d'architecture des ordinateurs (+mini-assembleur)}\nowrite
  \item {Éléments de systèmes d'exploitation}
\item Programmation structurée impérative (éléments de langage C)
\begin{itemize}
  \item {Structure d'un programme C}
  \item {Variables : déclaration (et initialisation), affectaction}
  \item {Évaluation d'expressions}
  \item {Instructions de contrôle : if, for, while}
  \item \alert{Types de données : entiers, caractères,
    réels}, {tableaux}, enregistrements
  \item \alert{Fonctions d'entrées/sorties (scanf/printf)}
  \item Écriture et appel de fonctions
  \item Débogage
\end{itemize}
\item {Notions de compilation
  \begin{itemize}
  \item Analyse lexicale, analyse syntaxique, analyse sémantique
  \item préprocesseur du compilateur C (include, define)
   \item Édition de lien
  \end{itemize}}
\item Algorithmes élémentaires
\item Méthodologie de résolution, manipulation sous linux
\end{itemize}
\end{frame}

% \section[Autres cours]{Survol des cours d'informatique des autres semestres}

% \begin{frame}
%   \frametitle{Survol des cours d'informatique des autres semestres}
% \pause
% \begin{itemize}
%     \item 
% Deuxième semestre\pause
%       \begin{itemize}
%         \item Programmation impérative (suite de EI)
%         \item Mini-projet, réalisation collaborative \hfill Sauf PC, SC.
%         \item Init. aux interfaces graphiques et au web \hfill Opt. maths, info.
%         \item Traitement automatique des langues \hfill Opt. maths, info.\pause
%       \end{itemize}
%     \item 
% Troisième semestre\pause
%       \begin{itemize}
%       \item Architecture et système \hfill Info, SPI.
%       \item Logique \hfill Info, (opt. maths)
%       \item Administration de parc info. \hfill Info, opt. PC, (opt. maths)
%       \item Calcul formel et prog. pour les sci. phy. \hfill PC.
%       \item Traitement de données app. à la finance \hfill MIEF, opt. maths.
%       \item Algorithmique numérique \hfill MIEF, opt. maths.
%      \end{itemize}
%     \item
% Quatrième semestre\pause
%       \begin{itemize}
%       \item Algorithmique et arbres \hfill Sauf PC, SC, SPI.
%       \item Programmation orientée objets \hfill Info.
%       \item Génie logiciel \hfill Info.
%       \item Systèmes réseaux \hfill Info, SPI.
%      \end{itemize}
%   \end{itemize}
% \end{frame}

\section[Représentation des données]{Représentation  des données}

\begin{frame}
  \frametitle{Représentation en binaire des données (rappel)\nowrite}
  \pause

  \begin{definition}[bit]
    \begin{itemize}
    \item Le chiffre binaire, ou \emph{bit}, est l'équivalent binaire
      de nos chiffres décimaux. Il peut valoir soit $0$ soit $1$. Un
      bit est une quantité élémentaire d'information (oui ou non,
      ouvert ou fermé, etc.).\pause
    \item L'information manipulée par un ordinateur est consituée de
      bits.
    \end{itemize}
  \end{definition}
\end{frame}

\subsection{Représentation des entiers naturels}
\begin{frame}
%3 minutes de copie pure compter 4.
 \frametitle{Représentation des entiers naturels}
  \begin{itemize}
  \item Dans une base donnée, un nombre entier positif est représenté
    de manière unique par une suite de chiffres de la base :\pause
    \begin{itemize}
    \item En base $10$, on écrit le nombre $109$ comme la suite des
      chiffres $1$, $0$, $9$ car :
\[
109 = 1\times 10^2 + 0\times 10^1 + 9 \times 10^0
\]
\pause
\item Et en base $2$, le nombre $25$ s'écrit comme la suite des bits
  $1$, $1$, $0$, $0$, $1$ car :
\[
\binaire{11001} = 1\times
      2^4 + 1\times 2^3 + 0\times 2^2 + 0\times 2^1 + 1\times 2^0
\]
\pause
\item Plus généralement, la correspondance entre la représentation et le nombre est donnée par :
\[
\binaire{b_k\ldots b_0} = \sum^{k}_{i = 0} b_i \times 2^i
\]
\end{itemize}
\end{itemize}
\end{frame}

\begin{frame}
  \frametitle{Les blagues d'informaticiens (ici XKCD)}
  \begin{center}
    \includegraphics[scale=0.5]{img/1_to_10.png}
  \end{center}
\pause
\begin{itemize}
\item
  Sur une échelle de $1$ à $10$, combien de chances y a t'il que cette
  question utilise le binaire ?\pause
\item $\ldots4$ ?\pause
\item Qu'est-ce qu'un $4$ ?
\end{itemize}
\end{frame}

\begin{frame}
  \begin{itemize}
\item<1-> $1$ bit représente $2$ possibilités, $2$ bits $4$ possibilités,
  $3$ bits $8$ possibilités, \ldots, $n$ bits $2^n$ possiblités.
\item<2-> Avec $n$ bits on peut écrire les $2^{n}$ premiers nombres de
  $0$ à $\sum^{n - 1}_{i = 0} 2^i = 2^n - 1$. En général, sur ordinateur, $n$ est
  fixé.
\item<3-> L'addition se fait comme en base 10, c'est même encore plus facile :
\uncover<4->{
  \begin{gather*}
    \begin{tabular}{l r l r}
       &\binaire{11001} & (= &25)\\ 
      + &\binaire{11} & (= & 3)\\ \hline \pause
      & \binaire{11100} & (= & 28)
    \end{tabular}
  \end{gather*}
}
\item<5-> Dans une représentation de taille fixée (sur ordinateur), il y a un risque de \alert{débordement} :
  \begin{gather*}
    \begin{tabular}{l r l r}
      &\binaire{11001} & (= & 25) \\
      + &\binaire{111} & (= &7)\\ \hline \pause
       & \alert{\binaire{00000}} & (= &\alert{0})
    \end{tabular}\qquad\uncover<6->{
   \begin{tabular}{l r l r}
      &\binaire{11011} & (= & 27) \\
      + &\binaire{111} & (= &7)\\ \hline 
       & \alert{\binaire{00010}} & (= &\alert{2})
    \end{tabular}}
  \end{gather*}
  \end{itemize}
\end{frame}


\subsection{Représentation des entiers relatifs}
\begin{frame}
 \frametitle{Représentation des entiers relatifs}
  \begin{itemize}
  \item Pour représenter les entiers relatifs, on peut réserver $1$
    bit au stockage du signe de l'entier ($0$ positif, $1$ négatif) et
    représenter en base $2$ sa valeur absolue sur les bits restants.\pause
\item Ça se fait parfois, mais il y a deux inconvénients :
    \begin{itemize}
      \item Il y a deux représentation du zéro \binaire{\alert{1}0000} et \binaire{\alert{0}0000}
      \item L'addition doit être modifiée\pause
    \end{itemize}
\item Alternative : la représentation en \alert{\emph{complément à deux}}. 
  \begin{itemize}
    \item Pour coder $-5$, on commence par coder $5$ : \hfill $\binaire{00101}$\pause 
\item On inverse les bits  \hfill $\binaire{11010}$\pause
\item On ajoute $1$ \hfill $\binaire{11011}$\pause
\item On obtient alors $-5$. L'addition avec $5$ fait zéro !
  \end{itemize}\pause
\item Ça fonctionne bien car :
  \begin{itemize}
    \item
  ajouter un nombre binaire avec le même nombre dont les bits ont
  étés inversés, donne \binaire{11\ldots1}.\pause
\item À nombre de bits fixé, ajouter $1$ donne zéro,
  par débordement.
  \end{itemize}\pause
\item Avec $n$ bits on représente les nombres de $-2^{n-1}$ à $2^{n-1} - 1$.
\item Remarque : le premier bit reste un bit de signe.
  \end{itemize}
\end{frame}

\subsection{Représentation des réels en virgule flottante}
\begin{frame}
    \frametitle{Représentation des réels en virgule flottante}
\pause
    \begin{itemize}
    \item La représentation informatique usuelle des réels s'inspire de la
      notation scientifique :
      \begin{align*}
        \pi &= 3,141592653589793\tag{pi}\\
        -700 \text{ milliards} &= -7 \times 10^{11}\tag{Paulson}\\
        h &= 6,626068 \times 10^{-34} \tag{Planck}\\
        \text{Univers} &= 1 \times 10^{80}  \tag{Atomes}
      \end{align*}\pause
\vspace{-.8cm}
\item Les bits sont séparés en :
  \begin{itemize}
  \item bit de signe\hfill (1 bit)
  \item mantisse \hfill (53 bits)
  \item exposant \hfill (11 bits)
  \end{itemize}\pause
\item En double précision (64 bits) :
  \begin{itemize}
\item exposant : entre $10^{-308}$ et $10^{308}$ (environ).
\item mantisse : 16 chiffres décimaux (environ).
  \end{itemize}\pause
\item Infini positif, infini négatif.
\item NaN : not a number.
    \end{itemize}
\end{frame}

\section[Types et E/S]{Types en C et entrées sorties associées}
\begin{frame}
  \frametitle{Types en C et entrées/sorties associées}
  
\pause
  \begin{itemize}
  \item Type des entiers relatifs  \alert{\C{int}} : 
    \begin{itemize}
    \item Déclaration et initialisation : \C{int n = -23;}.
    \item Représentation en complément à deux.
    \item E/S : \alert{\C{\%d}}.
    \end{itemize}
\pause
  \item Type des caractères \alert{\C{char}} :
    \begin{itemize}
    \item Déclaration et initialisation : \C{char c = 'A';}.
    \item Représentation sur 8 bits, ASCII, ISO-8859-x, UTF-8.
    \item E/S : \alert{\C{\%c}}.
    \end{itemize}
\pause
  \item Type des réels \alert{\C{double}} :
    \begin{itemize}
    \item Déclaration et initialisation : \C{double x = 3.14e-3;}.
    \item Représentation en virgule flottante sur 64 bits.
    \item E/S : \alert{\C{\%lg}} (mais plutôt \C{\%g} avec \C{printf}).
    \item \alert{Attention :} toujours mettre le point (équivalent
      anglais de la virgule) pour les constantes réelles (\C{1.0}).
    \end{itemize}
  \end{itemize}
\end{frame}

\begin{frame}
  \begin{block}{Entiers}
\C{int n;\\
...\\
printf("Entrer un nombre entier$\backslash$n");\\ \pause
scanf("\%d", \alert{\&}n);}
  \end{block}\pause

\begin{block}{Réels}
\C{double x;\\
...\\
printf("Entrer un nombre reel$\backslash$n");\\
scanf("\%lg", \&x);\\ \pause
printf("Vous avez saisi : \%g$\backslash$n", x);
}
\end{block}

\end{frame}

\begin{frame}
  
  \begin{block}{Caractères}
\C{char c;\\
...\\
printf("Entrer un caractère$\backslash$n");\\
scanf("\%c", \&c);
}
\\ \pause
    \alert{Attention :} mieux vaut utiliser \C{scanf(" \%c", \&c);}
    (voir démo)
  \end{block}\pause

%\structure{À emporter (pas au programme du premier partiel)\aemporter}
\begin{block}{Chaînes de caractères}
\C{char nom[64];\\
...\\
printf("Entrer votre nom$\backslash$n");\\
scanf("\%s", nom);\\ \pause
printf("Vous avez saisi : \%s$\backslash$n", nom);
}
\end{block}

\end{frame}

\pause
\section[Conversions]{Conversions automatiques entre types}
\begin{frame}
  \frametitle{Conversions  automatiques entre types}
  \begin{itemize}
\item Sans changement de représentation  :
  \begin{itemize}
    \item \C{char} vers \C{int}\pause
    \item \C{int} vers \C{char} (troncature)
  \end{itemize}\pause
  \item Avec changement de représentation :
    \begin{itemize}
    \item char ou entiers vers réels
    \item réels vers entiers ou char
    \end{itemize}
  \end{itemize}
\end{frame}


\section[Démos]{Démos et fin}
\begin{frame}
  \frametitle{Démos et fin}
\end{frame}
\end{document}
